\documentclass[12pt]{article}
\usepackage{amsmath, amssymb, geometry}
\geometry{margin=1in}
\setlength{\parindent}{0pt}
\setlength{\parskip}{0.6em}

\begin{document}

\begin{center}
\Large \textbf{Quiz 3\\ \vspace{0.6em}ECON 300 -- Labour Economics}\\

\vspace{0.6em}
\normalsize \textbf{March 04, 2026}
\end{center}

\vspace{0.5em}



\textbf{Part A: Multiple Choice Questions (8 marks)}

\textbf{1.} What is the primary difference between a monopsonist and a firm which operates in a perfectly competitive labour market?

A) The monopsonist is likely to be more profitable.

B) In order to recruit another worker, the monopsonist incurs a larger increase in labour costs.

C) In order to recruit another worker, the monopsonist incurs a smaller increase in labour costs.

D) The monopsonist hires such that the marginal revenue product of labour is equal to the marginal labour cost, but the non-monopsonist does not.

E) The labour supply curve is irrelevant for the monopsonist.

\vspace{1em}

\textbf{2.} Which of the following statements is false?

A) The monopsonist is a wage taker.

B) Technically speaking, the monopsonist has no demand curve for labour.

C) The monopsonist has market power in the labour market.

D) The monopsonist is a profit-maximizing firm.

E) The monopsonist may be able to carry out wage discrimination.

\vspace{1em}

\textbf{3.} Which of the following statements regarding monopsony is false?

A) The equilibrium wage and the equilibrium employment levels are lower than what they would be if the labour market were perfectly competitive.

B) All workers except for the marginal worker are paid more than their reservation wage.

C) If a competitive labour market were to become a monopsony, economic rent would be transferred from the workers directly to the employer.

D) The total wage bill is lower in the case of monopsony.

E) The marginal cost of labour is lower than the average cost of labour in the case of monopsony.

\newpage


\textbf{4.} If a firm operates in a perfectly competitive labour market, then:

A) It must be a monopsonist.

B) It faces an infinitely elastic supply of labour.

C) Its demand for labour must be inelastic.

D) It must also operate in a perfectly competitive output market.

E) The workers are unionized.

\vspace{0.8em}

\hrule
\vspace{0.6em}

\textbf{Part B: Problem (12 marks)}

Suppose the market labour supply and demand are:

\[
N^S = 25W
\]

\[
N^D = 500 - 12.5W
\]

where $W$ is the wage rate and $N$ is employment.

\begin{enumerate}

\item[(a)] (4 marks) Find the competitive equilibrium wage $W^*$ and employment $N^*$.

\item[(b)] (4 marks) Compute worker surplus (WS) and firm surplus (FS) at the competitive equilibrium. 

\textit{Hint:} Use geometric triangle areas. You may find it helpful to write inverse supply and inverse demand.

\item[(c)] (4 marks) Suppose labour demand increases to:

\[
N^{D'} = 600 - 12.5W
\]

Find the new equilibrium wage and employment. Briefly explain why both $W$ and $N$ change.

\end{enumerate}

\end{document}